\documentclass{theoremnote}

\setauthors{Ravi Dayabhai \& Tyler Campana}
\settitletext{Chord-cut Circles are $2$-colorable}
\setdate{May 2025}

\begin{document}

\begin{ProblemStatement}
  Prove that a circle cut by $n \in \mathbb{N}$ chords is $2$-colorable, i.e., the interior of a circle can be colored with at most $2$ colors such that no two regions sharing an edge are the same color.
\end{ProblemStatement}

\begin{ProofSolution}

  \begin{proof}

    We proceed by induction on the number of chords, $n$.

    \begin{induction}
    \basecase{
      When $n = 0$, the circle is not cut by any chords, so we can color the entire circle with one color. (Thus, it is trivially $2$-colorable.) Similarly, if $n = 1$, we can color the two sub-regions formed by the chord with $2$ different colors.
    }
    \inductivestep{
      Let $P(n)$ be the inductive hypothesis: a circle cut by $n$ chords is $2$-colorable. We want to show, assuming $P(n)$, that $P(n+1)$ holds, i.e., a circle cut by $n+1$ chords is also $2$-colorable.

      Assume that $P(n)$ is true -- if we start with a circle cut by $n$ chords, we can color the regions formed by these chords with no more than $2$ colors.
      Now, we add one more chord to the circle, so the circle is now cut by $n + 1$ chords. This chord will intersect some of the existing regions, so each region cut by the new chord will be divided into two sub-regions; however, now, these neighboring sub-regions share the same color.
      Next, to maintain the condition that no two adjacent regions can have the same color (the ``invariant''), we \textbf{recolor all regions on one side of the new chord with their opposite color}. 
      We can ensure that the invariant is not violated because:
      \begin{itemize}
        \item Every region \underline{on the unchanged side} of the new chord still has the same color as before, and this, original coloring satisfied the invariant (by the inductive hypothesis, $P(n)$).
        \item \underline{On the recolored side} of the new chord: 
        \begin{enumerate}
          \item \label{case:cut} For each region split into two sub-regions by the new chord, the resulting sub-region on the recolored side is colored with the opposite color, maintaining the invariant along the new chord. 
          \item The regions that were unaffected by the additional chord also satisfy the invariant because, while each switched colors, they respect the invariant among themselves. Also, the only recolored sub-regions they can possibly neighbor (from (\ref{case:cut})) also switched colors, so this side still satisfies the invariant. 
        \end{enumerate}
      \end{itemize}
    }
    We have shown that if $P(n)$ is true, then $P(n+1)$ is also true.
    \end{induction}

    Therefore, by induction, we conclude that a circle cut by $n$ chords is $2$-colorable.
  \end{proof}

\end{ProofSolution}

\end{document}
